\documentclass[conference]{IEEEtran}
\IEEEoverridecommandlockouts

\usepackage{cite}
\usepackage{amsmath,amssymb,amsfonts}
\usepackage{graphicx}
\usepackage{textcomp}
\usepackage{xcolor}
\usepackage{booktabs}
\usepackage{multirow}
\usepackage{microtype}

\def\BibTeX{{\rm B\kern-.05em{\sc i\kern-.025em b}\kern-.08em
    T\kern-.1667em\lower.7ex\hbox{E}\kern-.125emX}}

\begin{document}

\title{MinImageScorer: Error-Driven Difficulty Scoring\\
for Guided Copy-Paste Augmentation in\\
Agricultural Object Detection}


\author{
\IEEEauthorblockN{[FILL: Author Name\textsuperscript{1},
                  Supervisor Name\textsuperscript{1}]}
\IEEEauthorblockA{%
[FILL: Department of XXX, University of XXX, Country]}
}

\maketitle

% ----------------------------------------------------------
\begin{abstract}
We propose \textbf{MinImageScorer} (MIS), an error-driven framework
that scores per-object and image difficulty by combining model
confidence and localization penalties on training-set failure signals.
Validated on MinneApple and Global Wheat Head (GWHD), MIS reliably
ranks objects and images by difficulty and
correctly identifies domain- and size-level failure patterns
without requiring any addintional infomation.
When used to guide copy-paste augmentation, MIS-selected objects
outperform random selection on MinneApple ($\Delta\mathrm{AP}{=}
{+}0.013$); however, both strategies fall below baseline.
On GWHD, domain shift renders copy-paste ineffective regardless of
selection strategy.
These findings establish MIS as a reliable diagnostic for structural
difficulty in agricultural detection, while also revealing the
conditions under which copy-paste augmentation is unlikely to help.
\end{abstract}

\begin{IEEEkeywords}
object detection, copy-paste augmentation, difficulty scoring,
agricultural detection, small objects, domain shift
\end{IEEEkeywords}

% ----------------------------------------------------------
\section{Introduction}

Agricultural object detection benchmarks pose structural challenges
that are qualitatively different from general-purpose datasets such
as MS COCO. MinneApple~\cite{hani2020minneapple} is dominated by
densely packed, sub-resolution apple instances where detection
failure is driven primarily by scale rather than appearance.
Global Wheat Head Detection (GWHD)~\cite{david2021global} aggregates
images from seven geographic domains, introducing severe cross-domain
visual variation that causes systematic performance drops on
minority-domain images.
These two failure modes---small-object resolution limits and
cross-domain distributional shift---are well-documented by the
respective dataset authors, yet no prior augmentation method
explicitly diagnoses or targets them all at once automatically.

Copy-paste augmentation has demonstrated small gains on MS
COCO~\cite{kisantal2019augmentation, ghiasi2021simple}, where
objects span a broad size distribution and images are largely
within-distribution.
Applied blindly to structurally constrained datasets, however,
copy-paste risks compounding the very failure modes it aims to
address: pasting sub-resolution objects amplifies scale noise,
and pasting objects across domains amplifies distributional
mismatch.
No prior work provides an empirical criterion for assessing whether
copy-paste is appropriate for a given dataset \emph{before}
augmentation is applied.

We address this gap with two contributions.
First, we introduce \textbf{MinImageScorer} (MIS), a per-object
difficulty score derived directly from model failure signals.
We validate that MIS correctly identify dataset structural difficulty
without access to any additional information.
Second, we conduct a controlled study showing that score-guided
copy-paste outperforms random selection on the 
MinneApple dataset, while both strategies degrade below baseline. 
On the domain-shifted GWHD dataset, copy-paste
is ineffective regardless of selection strategy.
Taken together, these results show that MIS correctly identifies
where a model fails, and that the failure mode of a dataset itself---not the
selection strategy---determines whether copy-paste can succeed.

% ----------------------------------------------------------
\section{Methodology}

\subsection{Scoring Formulation}

Let $G(I)$ be the set of ground-truth objects in image $I$ and
$\mathcal{P}(g) = \{p : \mathrm{IoU}(p,g) \ge \tau\}$
the set of predictions matched to ground-truth object $g$.
Denoting $c_p = \mathrm{conf}(p)$ and $u_p = \mathrm{IoU}(p,g)$,
the \emph{Object Difficulty Score} is defined as:

\begin{equation}
S_{\mathrm{obj}}(g) =
\begin{cases}
  \min_{p \in \mathcal{P}(g)}
    \bigl[\alpha(1{-}c_p) + \beta(1{-}u_p)\bigr],
    & \mathcal{P}(g) \neq \varnothing, \\[3pt]
  \alpha + \beta, & \text{otherwise,}
\end{cases}
\label{eq:obj}
\end{equation}

where $\alpha$ and $\beta$ weight the confidence and localization
penalties, respectively.
% The $\min$ operator selects the best-matching prediction: an object
% is considered hard only if even its best-matched prediction is
% uncertain or poorly localized.
Missed objects (no prediction with $\mathrm{IoU} \ge \tau$) receive
the maximum score $\alpha{+}\beta$.
The \emph{Image Difficulty Score} combines mean object difficulty with
the false-positive rate:

\begin{equation}
S_{\mathrm{img}}(I)
  = w_1 \cdot \frac{1}{|G(I)|}
    \sum_{g\in G(I)} S_{\mathrm{obj}}(g)
  \;+\; w_2 \cdot \mathrm{FPrate}(I).
\label{eq:img}
\end{equation}

Object scoring uses a low confidence threshold ($\tau{=}0.01$)
to capture near-miss predictions; FP rate is computed at the
optimal operating threshold to reflect realistic deployment errors.
Weights $\alpha, \beta, w_1, w_2$ are calibrated by maximizing
Pearson correlation between $S_{\mathrm{img}}$ and per-image miss
rate and FP rate on the training set, requiring no domain-specific
prior knowledge.

\subsection{Score-Guided Copy-Paste Pipeline}

A baseline YOLOv8s detector is first trained on the target dataset.
Inference on the full training set yields per-prediction confidence
and IoU values, from which $S_{\mathrm{obj}}$ and $S_{\mathrm{img}}$
are computed for every annotated object and image.

Following the same pipeline as the author of \cite{kisantal2019augmentation} 
have done we will extend the dataset size 
by using Copy and Paste pipeline.
The special of our implementation is that instead of 
chosing a number of random objects in a random image and copy-paste it
multiple times in random locations with the same image,
we will use the score to change the probability of selecting an object 
and a image to copy and paste.

Under our \emph{score-guided copy-paste}, the probablity to choose and 
an ojbect and a background image are sampled
proportionally to their difficulty score via exponential
weighting---objects with high $S_{\mathrm{obj}}$ are drawn
more frequently, deliberately concentrating augmentation effort on
instances the baseline model most consistently fails to detect or
localize.
This new data then used to retrain a fresh YOLOv8s model
under identical hyperparameters, ensuring that any performance
difference is attributable to selection strategy alone.

% ----------------------------------------------------------
\section{Score Validation}
\label{sec:validation}

\subsection{Correlation with Detection Failures}

We can see in Table~\ref{tab:corr} that
all correlations are strongly positive, confirming
that the score reliably captures image-level difficulty.

\begin{table}[t]
\centering
\caption{Pearson correlation of $S_{\mathrm{img}}$ with per-image Miss rate and False
Positive rate}
\label{tab:corr}
\resizebox{\columnwidth}{!}{%
\begin{tabular}{@{}lccc@{}}
\toprule
\textbf{Dataset}
  & $r$(\,score,\,miss\,)
  & $r$(\,score,\,FP\,)
\\
\midrule
MinneApple  & 0.698 & 0.688 \\
GWHD        & 0.816 & 0.818 \\
\bottomrule
\end{tabular}%
}
\end{table}

\subsection{Domain and Size-Level Validity}

% The author of \cite{hani2020minneapple} identify
% small objects as the primary failure mode of MinneApple, while
% the author from \cite{david2021global} identify cross-domain visual
% shift as the defining challenge of GWHD.
% We test whether MIS recovers these patterns without access to
% any other information.

\noindent\textbf{Domain-level (GWHD).}
Table~\ref{tab:domain} show that the harder domains (lower AP) 
have higher Mean $S_{\mathrm{img}}$.
This indiate a strong relation between our scoring with GWHD failed mode.


\noindent\textbf{Size-level (MinneApple).}
Table~\ref{tab:size_score} shows that $S_{\mathrm{obj}}$ is
 higher for small objects than medium objects, which mean
our scoring has correctly
matched the known ap between these size categories.

Together, these results demonstrate that \textbf{MIS} can 
indentify the primary structural difficulty at both
domain and object-size granularity---using only training-set
model failure signals.

\begin{table}[t]
\centering
\caption{Correlation of mean $S_{\mathrm{img}}$ and Average Precision across GWHD domains}
\label{tab:domain}
\resizebox{\columnwidth}{!}{%
\begin{tabular}{@{}lcc@{}}
\toprule
\textbf{Domain}  & \textbf{AP}
  & \textbf{Mean} \boldmath{$S_{\mathrm{img}}$} \\
\midrule
rres\_1   & 0.572  & 0.200 \\
ethz\_1   & 0.550  & 0.227 \\
inrae\_1  & 0.523 & 0.251 \\
arvalis\_3 & 0.466 & 0.294 \\
usask\_1   & 0.408  & 0.312 \\
arvalis\_1 & 0.403 & 0.331 \\
arvalis\_2 & 0.365  & 0.424 \\
\bottomrule
\end{tabular}%
}
\end{table}

\begin{table}[t]
\centering
\caption{Correlation of Object Difficulty Score $S_{\mathrm{obj}}$ with object size on MinneApple.
\label{tab:size_score}}
\resizebox{\columnwidth}{!}{%
\begin{tabular}{@{}lccc@{}}
\toprule
\textbf{Dataset}
  & \textbf{Score Type}
  & \textbf{Small object (\boldmath{$\le$}32\,px\textbf)}
  & \textbf{Medium Object (\boldmath{$\le$}96\,px\textbf)} \\
\midrule
\multirow{2}{*}{MinneApple}
  & AP Score                        & 0.300 & 0.592  \\
  & Object Difficulty Score         & 0.378 & 0.222 \\

\bottomrule
\end{tabular}%
}
\end{table}



% ----------------------------------------------------------
\section{Augmentation Experiments}

\subsection{Setup}

We use YOLOv8s initialized from COCO pretrained weights.
We will use copy and paste augmentation to increase the dataset size by 30\%.
For each new image, 2--8 new objects are chosen to be pasted, applied
Scale jitter, random flipping following
Ghiasi~et~al.~\cite{ghiasi2021simple}.
The only difference between the two experimental conditions is
whether source objects are sampled uniformly (\textit{Random CP})
or proportionally to $S_{\mathrm{obj}}$ (\textit{Score CP}).

\subsection{Results}

\noindent\textbf{MinneApple.} Table~\ref{tab:ap} show that Score-guided 
copy-paste outperforms random selection across 
sub-metrics, confirming that
difficulty-weighted sampling directs augmentation toward
higher-signal objects.
However, both conditions fall below baseline indicates degraded
localization precision from pasting  objects without boundary
blending.  The AP$_S$ drop we believe reflects that score-selected
objects are predominantly sub-resolution and cannot be made more
learnable through repetition alone---a constraint imposed by the
dataset itself, not by the selection strategy.
% Backbone feature analysis (YOLOv8s P4, UMAP) confirms that
% augmented images are displaced from the test distribution,
% directly explaining the performance degradation.


\noindent\textbf{GWHD.}
Score-guided and random conditions are statistically
indistinguishable and both fall below baseline.
Score-guided sampling concentrates objects from domain-minority
sources, amplifying existing domain imbalance
rather than correcting it.
Feature-space analysis confirms that GWHD augmented images suffer
great distributional displacement, consistent with
the larger performance gap.
This we believe indicated that a domain-shifted dataset, 
selection quality is irrelevant if
the fundamental problem is covariate shift between augmented and
test images.

\begin{table}[t]
\centering
\caption{Copy-Paste Augmentation results}
\label{tab:ap}
\resizebox{\columnwidth}{!}{%
\begin{tabular}{@{}llrrrrr@{}}
\toprule
\textbf{Dataset} & \textbf{Condition}
  & \textbf{AP}
  & \textbf{AP\textsubscript{50}}
  & \textbf{AP\textsubscript{75}}
  & \textbf{AP\textsubscript{S}}
  & \textbf{AP\textsubscript{M}} \\
\midrule
\multirow{3}{*}{MinneApple}
  & Baseline          & 0.439 & 0.771 & 0.452 & 0.300 & 0.592 \\
  & Random CP         & 0.415 & 0.741 & 0.427 & 0.279 & 0.563 \\
  & \textbf{Score CP} & \textbf{0.428} & \textbf{0.759}
                      & \textbf{0.440} & \textbf{0.286}
                      & \textbf{0.581} \\
\midrule
\multirow{3}{*}{GWHD}
  & Baseline          & 0.507 & 0.922 & 0.491 & 0.136 & 0.501 \\
  & Random CP         & 0.486 & 0.905 & 0.456 & 0.118 & 0.480 \\
  & \textbf{Score CP} & \textbf{0.484} & \textbf{0.905}
                      & \textbf{0.455} & \textbf{0.098}
                      & \textbf{0.477} \\
\bottomrule
\end{tabular}%
}
\end{table}

% ----------------------------------------------------------
\section{Conclusion}

MinImageScorer reliably quantifies per-object detection difficulty
from model failure signals.

Score-guided copy-paste outperforms random selection on MinneApple
,confirming that difficulty-weighted sampling improves
selection quality.
Nevertheless, both strategies fall below baseline on both datasets
because the root failure modes---sub-resolution scale limits on
MinneApple and cross-domain distributional shift on GWHD---are
structural constraints that copy-paste cannot overcome when
augmented images are displaced from the test distribution,
as confirmed by backbone feature-space analysis.

The practical implication is precise: copy-paste is beneficial
when objects are underrepresented but visually compatible with the
test distribution; it is harmful when the failure mode is
structural (resolution limits or domain shift).
MIS provides a diagnostic criterion before augmentation: inspect
the size and domain distribution of high-score objects.
If they cluster in a regime the model cannot resolve regardless
of training exposure, copy-paste will not help.

Future work should explore two directions.
\textit{Better blending system}---Better blending method 
that can blend the boundary of the pasted object with the background
image to make the augmented image more natural and reduce the distributional shift.
\textit{Dynamic scoring}---recomputing object difficulty during
training rather than fixing scores upfront---would allow
augmentation to adapt as the model improves, avoiding continued
emphasis on objects that become learnable mid-training.
Both directions aim to keep augmented images within the test
distribution, addressing the core failure mode this work identifies.

% ----------------------------------------------------------
\begin{thebibliography}{00}

\bibitem{hani2020minneapple}
N.~Hani, P.~Roy, and B.~Bhanu,
``MinneApple: A Benchmark Dataset for Apple Detection and
Segmentation,''
\textit{IEEE Robotics and Automation Letters},
vol.~5, no.~2, pp.~852--858, 2020.

\bibitem{david2021global}
E.~David \textit{et~al.},
``Global Wheat Head Detection (GWHD) Dataset,''
\textit{Plant Phenomics}, vol.~2021, 2021.

\bibitem{kisantal2019augmentation}
M.~Kisantal \textit{et~al.},
``Augmentation for Small Object Detection,''
in \textit{Proc.\ ICACCI}, 2019.

\bibitem{ghiasi2021simple}
G.~Ghiasi \textit{et~al.},
``Simple Copy-Paste is a Strong Data Augmentation Method
for Instance Segmentation,''
in \textit{Proc.\ IEEE/CVF CVPR}, pp.~2918--2928, 2021.

\end{thebibliography}

\end{document}